%%%%%%%%%%%%%%%%%%%%%%%%%%%%%%%%%%%%%%%%%%%%%%%%%%%%%%%%%%%%%%%%%%%%%%%%
%%%%%%%%%%%%%%%%%%%%%% Simple LaTeX CV Template %%%%%%%%%%%%%%%%%%%%%%%%
%%%%%%%%%%%%%%%%%%%%%%%%%%%%%%%%%%%%%%%%%%%%%%%%%%%%%%%%%%%%%%%%%%%%%%%%

%%%%%%%%%%%%%%%%%%%%%%%%%%%%%%%%%%%%%%%%%%%%%%%%%%%%%%%%%%%%%%%%%%%%%%%%
%% NOTE: If you find that it says                                     %%
%%                                                                    %%
%%                           1 of ??                                  %%
%%                                                                    %%
%% at the bottom of your first page, this means that the AUX file     %%
%% was not available when you ran LaTeX on this source. Simply RERUN  %%
%% LaTeX to get the ``??'' replaced with the number of the last page  %%
%% of the document. The AUX file will be generated on the first run   %%
%% of LaTeX and used on the second run to fill in all of the          %%
%% references.                                                        %%
%%%%%%%%%%%%%%%%%%%%%%%%%%%%%%%%%%%%%%%%%%%%%%%%%%%%%%%%%%%%%%%%%%%%%%%%

%%%%%%%%%%%%%%%%%%%%%%%%%%%% Document Setup %%%%%%%%%%%%%%%%%%%%%%%%%%%%

% Don't like 10pt? Try 11pt or 12pt
\documentclass[10pt]{article}

% This is a helpful package that puts math inside length specifications
\usepackage{calc}
\usepackage{graphicx}
\usepackage{url}

%\usepackage{enumitem}

%\urlstyle{same}
% Simpler bibsection for CV sections
% (thanks to natbib for inspiration)
\makeatletter
\newlength{\bibhang}
\setlength{\bibhang}{1em}
\newlength{\bibsep}
 {\@listi \global\bibsep\itemsep \global\advance\bibsep by\parsep}
\newenvironment{bibsection}
    {\minipage[t]{\linewidth}\list{}{%
        \setlength{\leftmargin}{\bibhang}%
        \setlength{\itemindent}{-\leftmargin}%
        \setlength{\itemsep}{\bibsep}%
        \setlength{\parsep}{\z@}%
        }}
    {\endlist\endminipage}
\makeatother

% Layout: Puts the section titles on left side of page
\reversemarginpar

%
%         PAPER SIZE, PAGE NUMBER, AND DOCUMENT LAYOUT NOTES:
%
% The next \usepackage line changes the layout for CV style section
% headings as marginal notes. It also sets up the paper size as either
% letter or A4. By default, letter was used. If A4 paper is desired,
% comment out the letterpaper lines and uncomment the a4paper lines.
%
% As you can see, the margin widths and section title widths can be
% easily adjusted.
%
% ALSO: Notice that the includefoot option can be commented OUT in order
% to put the PAGE NUMBER *IN* the bottom margin. This will make the
% effective text area larger.
%
% IF YOU WISH TO REMOVE THE ``of LASTPAGE'' next to each page number,
% see the note about the +LP and -LP lines below. Comment out the +LP
% and uncomment the -LP.
%
% IF YOU WISH TO REMOVE PAGE NUMBERS, be sure that the includefoot line
% is uncommented and ALSO uncomment the \pagestyle{empty} a few lines
% below.
%

%% Use these lines for letter-sized paper
\usepackage[paper=letterpaper,
            %includefoot, % Uncomment to put page number above margin
            marginparwidth=1.2in,     % Length of section titles
            marginparsep=.05in,       % Space between titles and text
            margin=1in,               % 1 inch margins
            includemp]{geometry}

%% Use these lines for A4-sized paper
%\usepackage[paper=a4paper,
%            %includefoot, % Uncomment to put page number above margin
%            marginparwidth=30.5mm,    % Length of section titles
%            marginparsep=1.5mm,       % Space between titles and text
%            margin=25mm,              % 25mm margins
%            includemp]{geometry}

%% More layout: Get rid of indenting throughout entire document
\setlength{\parindent}{0in}

%% This gives us fun enumeration environments. compactitem will be nice.
\usepackage{paralist}

%% Reference the last page in the page number
%
% NOTE: comment the +LP line and uncomment the -LP line to have page
%       numbers without the ``of ##'' last page reference)
%
% NOTE: uncomment the \pagestyle{empty} line to get rid of all page
%       numbers (make sure includefoot is commented out above)
%
\usepackage{fancyhdr,lastpage}
\pagestyle{fancy}
%\pagestyle{empty}      % Uncomment this to get rid of page numbers
\fancyhf{}\renewcommand{\headrulewidth}{0pt}
\fancyfootoffset{\marginparsep+\marginparwidth}
\newlength{\footpageshift}
\setlength{\footpageshift}
          {0.5\textwidth+0.5\marginparsep+0.5\marginparwidth-2in}
\lfoot{\hspace{\footpageshift}%
       \parbox{4in}{\, \hfill %
                    \arabic{page} of \protect\pageref*{LastPage} % +LP
%                    \arabic{page}                               % -LP
                    \hfill \,}}

% Finally, give us PDF bookmarks
\usepackage{color,hyperref}
\definecolor{darkblue}{rgb}{0.0,0.0,0.3}
\hypersetup{colorlinks,breaklinks,
            linkcolor=darkblue,urlcolor=darkblue,
            anchorcolor=darkblue,citecolor=darkblue}

%%%%%%%%%%%%%%%%%%%%%%%% End Document Setup %%%%%%%%%%%%%%%%%%%%%%%%%%%%


%%%%%%%%%%%%%%%%%%%%%%%%%%% Helper Commands %%%%%%%%%%%%%%%%%%%%%%%%%%%%

% The title (name) with a horizontal rule under it
%
% Usage: \makeheading{name}
%
% Place at top of document. It should be the first thing.
\newcommand{\makeheading}[1]%
        {\hspace*{-\marginparsep minus \marginparwidth}%
         \begin{minipage}[t]{\textwidth+\marginparwidth+\marginparsep}%
                {\large \bfseries #1}\\[-0.15\baselineskip]%
                 \rule{\columnwidth}{1pt}%
         \end{minipage}}

% The section headings
%
% Usage: \section{section name}
%
% Follow this section IMMEDIATELY with the first line of the section
% text. Do not put whitespace in between. That is, do this:
%
%       \section{My Information}
%       Here is my information.
%
% and NOT this:
%
%       \section{My Information}
%
%       Here is my information.
%
% Otherwise the top of the section header will not line up with the top
% of the section. Of course, using a single comment character (%) on
% empty lines allows for the function of the first example with the
% readability of the second example.
\renewcommand{\section}[2]%
        {\pagebreak[2]\vspace{1.3\baselineskip}%
         \phantomsection\addcontentsline{toc}{section}{#1}%
         \hspace{0in}%
         \marginpar{
         \raggedright \scshape #1}#2}

% An itemize-style list with lots of space between items
\newenvironment{outerlist}[1][\enskip\textbullet]%
        {\begin{itemize}[#1]}{\end{itemize}%
         \vspace{-.6\baselineskip}}

% An environment IDENTICAL to outerlist that has better pre-list spacing
% when used as the first thing in a \section
\newenvironment{lonelist}[1][\enskip\textbullet]%
        {\vspace{-\baselineskip}\begin{list}{#1}{%
        \setlength{\partopsep}{0pt}%
        \setlength{\topsep}{0pt}}}
        {\end{list}\vspace{-.6\baselineskip}}

% An itemize-style list with little space between items
\newenvironment{innerlist}[1][\enskip\textbullet]%
        {\begin{compactitem}[#1]}{\end{compactitem}}

% An environment IDENTICAL to innerlist that has better pre-list spacing
% when used as the first thing in a \section
\newenvironment{loneinnerlist}[1][\enskip\textbullet]%
        {\vspace{-\baselineskip}\begin{compactitem}[#1]}
        {\end{compactitem}\vspace{-.6\baselineskip}}

% To add some paragraph space between lines.
% This also tells LaTeX to preferably break a page on one of these gaps
% if there is a needed pagebreak nearby.
\newcommand{\blankline}{\quad\pagebreak[2]}

% Uses hyperref to link DOI
\newcommand\doilink[1]{\href{http://dx.doi.org/#1}{#1}}
\newcommand\doi[1]{doi:\doilink{#1}}


%%%%%%%%%%%%%%%%%%%%%%%% End Helper Commands %%%%%%%%%%%%%%%%%%%%%%%%%%%

%%%%%%%%%%%%%%%%%%%%%%%%% Begin CV Document %%%%%%%%%%%%%%%%%%%%%%%%%%%%


\begin{document}

\large

\makeheading{R\&D Projects and Publications}



\section{R\&D Projects (2013 - 2020)} \begin{bibsection}
     \item An artificial Intelligence based robot for sorting plastic bottles: \newline It is a joint work with a trash recycling company. The objective is to detect, classify and separate plastic bottles from other garbage placed on a moving conveyer belt. For detection and classification, we used hyper-spectral camera and for bottle collection, we used an Ethercat based manipulator. I worked on the development of image processing and machine learning algorithms. Also, I worked on low level system integration and designed the whole system architecture.   
     
	\item A mobile platform for SLAM with multiple cameras, an IMU, and a GPS: \newline It is a joint work with a Researcher Center. The main system has been developed using ROS with C++ and Python. I worked as the system architect. I developed the platform software. Also, I contributed to the implementation of SLAM algorithms. 
	
    \item A Industry 4.0 based project: \newline It is a joint work with a large manufacturing company. We worked on \emph{``Big Data''} based predictive maintenance of various parts in the main production line.  
    
    \item TRUST - a humanoid robot: \newline It is a humanoid robot with 35 major joints capable of doing biped walking, stair climbing, and simple manipulation. We used both Ethercat and CAN communication and Elmo motor controllers to control most of the joints.  For some of the joints, we used our own motor controllers. The overall motion system has been developed using C, C++ on Xenomai. I designed the whole software system architecture and implemented software platform. I also contributed to the controller design and implementations. 

\item Tibo - a Ski-robot: \newline We participated in "Winter Olympics 2018: robot ski tournament" with Tibo. I designed the overall system architecture and implemented the platform. Also I contributed to trajectory generation algorithms. 

\item RailBot: \newline It is a small rail robot for remote monitoring: This robot has an image camera and a thermal camera. It moves autonomously and transfers images and other information to some remote user. I worked on the development of the motor controller to control the speed of the robot. I also developed the mobile app for android devices. 

\item A temperature monitoring system: \newline It is a temperature monitoring system of cooling water for industrial use: We developed it for some manufacturing company. I developed the whole firmware on Microcontroller - STM32.  

\end{bibsection}

\section{Dissertations} \begin{bibsection}
     \item Mohammad Khairul Hasan
     \emph{Approximation Algorithms for Facility Location and Graph Partitioning
Problems}. (Ph.D.). Korea Advanced Institute of Science and Technology, 2013
     \item Mohammad Khairul Hasan
     \emph{Improved approximation algorithms for connected facility location problems}. (MS). Korea Advanced Institute of Science and Technology, 2006
\end{bibsection}

\section{Journal Publications} \begin{bibsection}
     \item Mohammad Khairul Hasan, and Kyung-Yong Chwa.
     \emph{Approximation algorithms for the Weighted t-Uniform Sparsest Cut and some other graph partitioning problems}.J. Comput. Syst. Sci. 82(6): 1044-1063, 2016.
	\item Md. Shamim Ahsan, Man Seop Lee, Mohammad Khairul Hasan, Young-Chul Noh, Farid Ahmed, Martin B.G. Jun.
     \emph{Formation mechanism of self-organized nano-ripples on quartz surface using femtosecond laser pulses
}.Optik - International Journal for Light and Electron, 126(24), 2015.     
    \item Kunho Kim, Mohammad K. Hasan, Jae-Pil Heo, Yu-Wing Tai, Sung-eui
    Yoon. \emph{Probabilistic cost model for nearest neighbor search in image
    retrieval} Computer Vision and Image Understanding, 116(9): 991-998, 2012.
    \item Hyunwoo Jung, Mohammad Khairul Hasan, and Kyung-Yong Chwa.
     \emph{A 6.5 factor primal-dual approximation algorithm for
        the connected facility location problem}.J. Comb. Optim., 18(3):258-271, 2009.

    \item Mohammad Khairul Hasan, Hyunwoo Jung, and Kyung-Yong Chwa. \emph {Approximation
        algorithms for connected facility location problems}. J. Comb. Optim., 16(2):155-172,
        2008.

\end{bibsection}


\section{Conference Publications} \begin{bibsection}

    \item Mohammad Khairul Hasan, Sung-Eui Yoon, and Kyung-Yong Chwa. \emph{Bounds on the
        geometric mean of arc lengths for bounded-degree planar graphs}.
        In FAW, pages 153-162, 2009.
    \item Hyunwoo Jung, Mohammad Khairul Hasan, and Kyung-Yong Chwa. \emph{Improved primal-dual approximation
        algorithm for the connected facility location problem}. In COCOA, pages 265-277, 2008.
    \item Mohammad Khairul Hasan, Hyunwoo Jung, and Kyung-Yong Chwa. \emph{Improved approximation
        algorithm for connected facility location problems}. In COCOA, pages 311-322, 2007.

\end{bibsection}


\end{document}

%%%%%%%%%%%%%%%%%%%%%%%%%% End CV Document %%%%%%%%%%%%%%%%%%%%%%%%%%%%%
